%% P01 --- Liczby zmiennoprzecinkowe: co maszyna naprawdę liczy
%%
%% Zalążek. Poniższy opis jest kontraktem tego programu: co ma uzasadnić i co
%% ma dać czytelnikowi. Napisanie programu oznacza usunięcie bloku
%% \programstub{} i zastąpienie go programem. Jeśli po przerobieniu materiału
%% opis okaże się błędny, popraw go i odnotuj sprzeczność w CLAUDE.md w sekcji
%% *Resolved questions*.

\program{Liczby zmiennoprzecinkowe: co maszyna naprawdę liczy}
\label{prog:P01}

\programstub{%
\textit{[Opis do przetłumaczenia --- patrz wydanie angielskie.]}

Argument: a float is a sign, an exponent and a fixed number of significant
bits, so the gap between representable numbers grows with magnitude, and
equality is not a question you should ask. Payoff: why \code{bf16} is the
training format and \code{fp16} needs loss scaling --- they have the same
width and different exponent budgets, and the reader computes the overflow
threshold rather than being told it; why a sum of gradients is not
associative and two GPUs therefore disagree in the last bits; what machine
epsilon is and why comparing two loss values that differ in the eighth
decimal is measuring the hardware.

\medskip
\textbf{Planowana długość:} 45 ramek.
}
